\section{Stable detection, its Galois connection, and its losses}

\subsection{The pointed spectrum attached to an unpointed map}

\begin{claim}{F.1}{Proved}{The spectrum in the proposed assignment}
For a map $f:X\to Y$ of $\mathcal U$-small spaces, put
\[
 C_f=Y\amalg_X^h *,\qquad
 D(f)=\cofib(\Sigma^\infty_+X\xrightarrow{\Sigma^\infty_+f}
                   \Sigma^\infty_+Y)\in\Sp.
\]
The cone point gives $C_f$ a specified basepoint. There is a natural equivalence
$D(f)\simeq\Sigma^\infty C_f$, with \emph{reduced} suspension spectrum on the
right. Moreover, $D(C_f\to *)\simeq\Sigma D(f)$. Both assertions include
empty and disconnected $X,Y$.
\end{claim}
\begin{proof}
Adjoining a disjoint basepoint is a left adjoint from spaces to pointed spaces.
It therefore sends the defining homotopy pushout for $C_f$ to the corresponding
pointed homotopy pushout, with a copy of $S^0$ at the cone vertex. Collapsing
that copy of $S^0$ yields the pointed cofiber of $f_+:X_+\to Y_+$.
Applying the colimit-preserving functor $\Sigma^\infty$ gives the asserted
equivalence. Equivalently, this is the pasting law for the two pushout squares
which first attach the cone and then collapse its vertex together with the
disjoint basepoint. For the second assertion the specified cone point splits
the augmentation $\Sigma^\infty_+C_f\to\mathbb S$. Its fiber is
$\Sigma^\infty C_f$; its cofiber is its suspension. For $X=\varnothing$ the
first formula reads $D(f)=\Sigma^\infty_+Y$ and $C_f=Y_+$, so that no
unmentioned nonemptiness assumption has entered the calculation.
\end{proof}

\begin{claim}{F.2}{Proved}{Localizing subcategories give reversible excision}
Let $\mathcal L$ be a full replete localizing subcategory of
$\Ho(\Sp_{\mathcal U})$: it is closed under all shifts, distinguished
triangles, retracts, and $\mathcal U$-small coproducts. Define
\[
 R(\mathcal L)=\{u:A\to B\text{ in }\Cat\mid D(Bu)\in\mathcal L\}.
\]
Then $R(\mathcal L)$ is a strongly saturated fundamental localizer in the
relative-Theorem-A convention of this paper. It satisfies Loc (6 a),
Loc (6 b), and Loc (6 c), and for every functor $u$ it satisfies
\[
 u\in R(\mathcal L)\quad\Longleftrightarrow\quad
 C_{Bu}\text{ is }R(\mathcal L)\text{-aspherical}.
\]
The asphericity of a space means that of any category-of-simplices model of
that space; the comparison weak equivalences make the definition independent
of the model.
\end{claim}
\begin{proof}
The class $\mathcal L$ is closed under homotopy colimits of spectra. Indeed,
sequential homotopy colimits are cofibers of the shift-minus-identity map on
a coproduct. Geometric realizations are sequential homotopy colimits of
their skeletal filtrations, whose successive cofibers are suspensions of
coproducts. The simplicial replacement expresses an arbitrary small homotopy
colimit in these terms. Cofiber commutes with homotopy colimits. It follows
that maps with cofiber in $\mathcal L$ are closed under homotopy colimits in
the arrow category. The octahedral triangle for composable maps gives two
out of three; functoriality of cofiber gives closure under arrow retracts.
The Verdier quotient by the thick subcategory $\mathcal L$ detects precisely
these maps: a map becomes invertible exactly when its cone belongs to
$\mathcal L$. The composite $\Cat\xrightarrow{B}\cS
\xrightarrow{\Sigma^\infty_+}\Sp\to\Ho(\Sp)/\mathcal L$ therefore
detects $R(\mathcal L)$ and proves its strong saturation. This quotient can
be taken in the larger universe; no smallness of its hom sets in
$\mathcal U$ is assumed.

If $A$ has a terminal object, then $BA$ is contractible, so $A\to e$ is in
$R(\mathcal L)$. To check the relative Theorem A, take a commutative triangle
$A\xrightarrow{u}B\to C$ and suppose $A/c\to B/c$ belongs to
$R(\mathcal L)$ for every $c\in C$. Postcomposition makes
$c\mapsto A/c$ and $c\mapsto B/c$ covariant diagrams on $C$.
There are natural weak equivalences
\[
 \hocolim_{c\in C}B(A/c)\longrightarrow BA,\qquad
 \hocolim_{c\in C}B(B/c)\longrightarrow BB.
\]
For clarity, the category-theoretic projection
$p:\int_C(A/c)\to A$ has left adjoint
$a\mapsto(v(a),a,\id_{v(a)})$, where $v:A\to C$ is the structure functor.
The adjunction supplies the weak equivalence after the usual nerve
comparison for the Grothendieck construction. The construction is natural
in $u$. Applying $\Sigma^\infty_+$, taking the cofiber of the comparison,
and using homotopy-colimit closure proves that $D(Bu)\in\mathcal L$.

For a homotopy pushout square of spaces with vertices $X_0,X_1,X_2,X_3$,
the suspension spectra form the distinguished triangle
\[
 \Sigma^\infty_+X_0\longrightarrow
 \Sigma^\infty_+X_1\oplus\Sigma^\infty_+X_2\longrightarrow
 \Sigma^\infty_+X_3\longrightarrow\Sigma\Sigma^\infty_+X_0.
\]
For a map of such squares, with vertical maps $f_i$, taking cofibers gives
the triangle $D(f_0)\to D(f_1)\oplus D(f_2)\to D(f_3)$.
The three choices of two known zero terms in the quotient prove respectively
Loc (6 a), (6 b), and (6 c); in the last case each summand is a retract of
the middle term. Mayer--Vietoris squares of sieve covers become homotopy
pushouts on nerves. The same proof applies to $R(\mathcal L)$-cocartesian
squares, since their defining comparison is inverted in the quotient.
Finally, \claimref{F.1} and closure under shifts identify the two sides of
the asserted cofiber criterion.
\end{proof}

\begin{claim}{F.3}{Proved}{Accessibility and the size of the assignments}
If $\mathcal L$ is generated by a $\mathcal U$-small set of spectra, then
$R(\mathcal L)$ is accessible. If a fundamental localizer $\W$ is generated
by a $\mathcal U$-small set $S$ of functors, then
\[
 T(\W):=\operatorname{loc}\{D(Bu):u\in\W\}
       =\operatorname{loc}\{D(Bs):s\in S\}.
\]
Without these set-generation hypotheses the order-theoretic assignments
still exist in the larger universe; we make no assertion that the
corresponding localizations are accessible in $\mathcal U$.
\end{claim}
\begin{proof}
A set-generated stable localizing subcategory has a presentable stable
quotient, obtained by localization at the set of all integer shifts of its
generators mapped to zero. The quotient functor $q:\Sp\to\Sp/\mathcal L$
is accessible and colimit preserving; existence can equivalently be
obtained by the transfinite construction in
\citet[Theorem 3.7]{FCasacubertaGutierrezRosicky2014}, applied to the
coproduct of the shifted generators. The inverse image of equivalences
under the accessible arrow functor induced by
$q\Sigma^\infty_+:\cS\to\Sp/\mathcal L$ is accessible: it is the
iso-comma category of this arrow functor and the accessible inclusion of
equivalences. Choose a regular cardinal for which this category and its
inclusion into the arrow category are accessible. A small collection of
representatives of its presentable arrows generates the whole class under
filtered colimits. The class is strongly saturated and closed under all
colimits, so these representatives generate it as a localization class.
Transport to $\Cat$ uses the nerve/category-of-simplices correspondence;
the accessibility statement is precisely the corresponding part of
\citet[Theorem 4.2.15]{RootCisinski2006}. For $\W$ generated by $S$, let
$\mathcal L_S=\operatorname{loc}\{D(Bs):s\in S\}$. By \claimref{F.2},
$R(\mathcal L_S)$ is fundamental and contains $S$, hence contains $\W$.
Thus $T(\W)\subseteq\mathcal L_S$. The reverse inclusion follows from
$S\subseteq\W$.
\end{proof}

\begin{claim}{F.4}{Proved}{The Galois connection}
Order fundamental localizers and full localizing subcategories of
$\Ho(\Sp)$ by inclusion. With $T$ and $R$ as above,
\[
 T(\W)\subseteq\mathcal L
   \quad\Longleftrightarrow\quad
 \W\subseteq R(\mathcal L).
\]
Consequently $RT$ is an increasing idempotent closure operator on
fundamental localizers, while $TR$ is a decreasing idempotent operator on
localizing subcategories. The assignments restrict to the accessible and
set-generated objects described in \claimref{F.3}.
\end{claim}
\begin{proof}
The left-hand inclusion holds exactly when $D(Bu)\in\mathcal L$ for every
$u\in\W$, which is the right-hand inclusion. Both operations preserve
inclusions. Applying the equivalence to the identities gives
$\W\subseteq RT(\W)$ and $TR(\mathcal L)\subseteq\mathcal L$.
Applying $R$ and $T$ to these two inequalities in turn gives
$RTR=R$ and $TRT=T$, and hence the two idempotence assertions.
\end{proof}

\begin{claim}{F.5}{Proved}{The exact image and the fibers of stable detection}
Let $\mathscr D$ denote the class of spectra equivalent to
$\Sigma^\infty K$ for a pointed $\mathcal U$-small space $K$.
Then
\[
 TR(\mathcal L)=\operatorname{loc}(\mathcal L\cap\mathscr D).
\]
The stable fixed objects are exactly the localizing subcategories generated
by the suspension spectra which they contain. The unstable fixed objects
are exactly the localizers $\W=RT(\W)$, equivalently those of the form
$R(\mathcal L)$. On these two fixed-object posets, $T$ and $R$ are inverse
order isomorphisms. Furthermore,
\[
 R(\mathcal L)=R(\mathcal M)
 \quad\Longleftrightarrow\quad
 \mathcal L\cap\mathscr D=\mathcal M\cap\mathscr D.
\]
Thus each nonempty fiber of $R$ has the least member $TR(\mathcal L)$;
its other members are exactly the localizing subcategories with the same
intersection with $\mathscr D$. Each such fiber is order convex.
\end{claim}
\begin{proof}
By \claimref{F.1}, every $D(Bu)$ belongs to $\mathscr D$. Conversely every
pointed space $K$ is, as a pointed space up to natural equivalence, the
homotopy cofiber of its basepoint map $*\to K$. Replace this map by the
functor on categories of simplices of its singular simplicial sets. The
natural realization comparisons are weak equivalences at its source and
target and hence on its cofiber. Therefore every spectrum in
$\mathcal L\cap\mathscr D$ occurs, up to equivalence, among the generators
of $TR(\mathcal L)$. This proves the formula. The fixed-object assertions
follow from \claimref{F.4}. The same realization argument shows that
equality of the two map classes tests every object of $\mathscr D$ and
therefore is equivalent to equality of the intersections. If
$\mathcal L_0\subseteq\mathcal L_1\subseteq\mathcal L_2$ and the two
outside members have the same intersection with $\mathscr D$, the middle
one has that intersection as well. This proves order convexity without
asserting that the fiber has a greatest member.
\end{proof}

\begin{claim}{F.6}{Proved}{The zero stable subcategory detects integral homology}
One has $R(0)=\WZ$. In fact
\[
 T(\W)=0\quad\Longleftrightarrow\quad\W\subseteq\WZ.
\]
In particular, every fundamental localizer between $\Winfty$ and $\WZ$
has the same image $0$ under $T$, and $RT(\W)=\WZ$ throughout this interval.
\end{claim}
\begin{proof}
The spectrum $D(Bu)$ is connective, since it is the cofiber of a map of
connective spectra. A connective spectrum $V$ with $H\mathbb Z\wedge V=0$
must be zero. Otherwise, choose its least nonzero homotopy group
$\pi_nV$. The stable Hurewicz theorem identifies this first nonzero
homotopy group with $H_n(V;\mathbb Z)$, a contradiction. Thus
\[
 D(Bu)=0\quad\Longleftrightarrow\quad
 H\mathbb Z\wedge D(Bu)=0
 \quad\Longleftrightarrow\quad u\in\WZ.
\]
The last equivalence is the exact homology sequence of the cofiber of
$\Sigma^\infty_+Bu$, including degree zero. A localizing subcategory
generated by a family is zero precisely when each member of that family
is zero, which proves the second assertion.
\end{proof}

\begin{claim}{F.7}{Refuted}{The proposed bijection with all localizing subcategories}
The assignments in Lead 6 are not inverse bijections, even with
accessibility and set generation. Let
\[
 \mathcal L_{\mathbb Z}
   =\{V\in\Ho(\Sp):H\mathbb Z\wedge V=0\}.
\]
Then $\mathcal L_{\mathbb Z}\ne0$, but
$R(\mathcal L_{\mathbb Z})=R(0)=\WZ$ and
$TR(\mathcal L_{\mathbb Z})=0$.
\end{claim}
\begin{proof}
The equality of the two map classes is the connective argument of
\claimref{F.6}. A nonzero Morava spectrum $K(n)$ at any prime, with
$n>0$, belongs to $\mathcal L_{\mathbb Z}$: one has
$H\mathbb Z\wedge K(n)=0$, while
$\pi_*K(n)=\mathbb F_p[v_n,v_n^{-1}]\ne0$.
This standard example is recorded explicitly in
\citet[Example 9.20]{FLawson2019}. Hence the two stable subcategories
are distinct. For an immediately singly generated counterexample, use
$\operatorname{loc}(K(n))\subseteq\mathcal L_{\mathbb Z}$ in place of
$\mathcal L_{\mathbb Z}$; it is nonzero and has exactly the same
intersection with connective spectra as $0$. The kernel of a
homological localization of spectra is also generated by a set;
equivalently Bousfield's smallness argument gives a
single coproduct of generators. The set-generation assertion is also
recalled in \citet[Section 2, discussion of Lemma 1.13 of Bousfield]{FHoveyPalmieri1999}.
Thus imposing accessibility does not remove this counterexample.
The loss occurs because the input cofibers are suspension spectra, and
in particular connective spectra; no space map probes a nonzero
$H\mathbb Z$-acyclic connective spectrum.
\end{proof}

\subsection{An explicit unstable distinction invisible to the stable assignment}

\begin{claim}{F.8}{Proved}{A homology equivalence between free-group categories}
Let $F=\langle x,y\rangle$ be the free group on two generators. The
one-object category $BF$ has arrows the reduced words in
$x,x^{-1},y,y^{-1}$, identity the empty word, and composition concatenation
followed by free reduction. The endomorphism
\[
 \alpha:F\longrightarrow F,\qquad
 \alpha(x)=x,\quad \alpha(y)=yxyx^{-1}y^{-1}
\]
defines a functor $B\alpha:BF\to BF$ which is an integral homology
equivalence, and indeed an $E$-homology equivalence for every spectrum
$E$, but is not a weak homotopy equivalence. The group $F$ has no
nontrivial perfect subgroup.
\end{claim}
\begin{proof}
The graph with one vertex and two oriented loops is a $K(F,1)$: its
universal covering graph is the tree of reduced words. Therefore
$H_0(BF;\mathbb Z)=\mathbb Z$, $H_1(BF;\mathbb Z)=\mathbb Z^2$, and
all higher groups vanish. Counting exponents in the two words shows that
$H_1(B\alpha;\mathbb Z)$ is the identity matrix. The functor is therefore
in $\WZ$, and \claimref{F.6} says that its spectral cofiber is zero.
Smashing this cofiber with any $E$ proves the generalized homology assertion.

To prove failure of surjectivity without relying on a word-search
heuristic, let $h_0,h_1:F\to S_3$ be the unique homomorphisms with
\[
 h_0(x)=h_1(x)=(12),\qquad
 h_0(y)=(12),\qquad h_1(y)=(13).
\]
They are distinct. The braid relation for the two transpositions gives
$h_0\alpha(y)=h_1\alpha(y)=(12)$, and equality on $x$ is immediate.
Thus $h_0\alpha=h_1\alpha$. Surjectivity of $\alpha$ would force
$h_0=h_1$, a contradiction. In particular the induced fundamental-group
map is not an isomorphism. The program
\texttt{code/f\_free\_group\_witness.py} checks all compositions and
all 216 associativity triples of the six-arrow target category $BS_3$,
the identity and inverse laws, both evaluations, and the exponent matrix;
its complete output is attached. The program supplements the universal
property proof; it does not pretend to enumerate the infinite category $BF$.

Finally, any subgroup of $F$ is the fundamental group of a connected
covering of the two-loop graph. Choosing a maximal tree in that covering
expresses the subgroup as a free group on its remaining edges. A nontrivial
free group has nonzero free abelianization, so no such subgroup is perfect.
\end{proof}

\subsection{Plus construction and acyclic nullification}

\begin{claim}{F.9}{Proved}{The plus localizer has the same aspherical objects as integral homology}
Let $P_+$ be nullification at the universal acyclic space of
\citet[Theorem 4.3]{FBerrickCasacuberta1999}, applied componentwise to all
spaces, with $P_+(\varnothing)=\varnothing$. It is the plus construction
with respect to the largest perfect subgroup of each component's fundamental
group. Put
$\W_+=\{u:A\to B\mid P_+(Bu)\text{ is a weak equivalence}\}$.
Then $\W_+$ is an accessible fundamental localizer, and
\[
 \Winfty\subsetneq\W_+\subsetneq\WZ.
\]
Its aspherical categories are precisely the nonempty connected integrally
acyclic categories, hence are exactly the $\WZ$-aspherical categories.
Nevertheless $\W_+$ fails Loc (6 b), and $T(\W_+)=0$.
\end{claim}
\begin{proof}
The cited theorem supplies a single small connected acyclic space whose
nullification realizes the indicated plus construction. Localization at one
space-to-point map is accessible. Its equivalences form a strongly saturated
class closed under homotopy colimits; the argument with the diagrams
$A/c$ in \claimref{F.2} proves the relative Theorem A after transfer to
categories. Thus it defines a fundamental localizer.

The plus-construction unit $X\to P_+X$ is an integral homology isomorphism
on each component. Naturality and two out of three imply
$\W_+\subseteq\WZ$. If $X$ is nonempty, connected, and integrally
acyclic, then $\pi_1X$ is perfect, so $P_+X$ is simply connected and
integrally acyclic. Hurewicz, in the least possible nonzero degree, and
Whitehead imply that $P_+X$ is weakly contractible. Conversely,
$P_+X\simeq *$ and the homology isomorphism of the unit imply that $X$
is connected and acyclic. The empty space is not made into a point by
this localization.

The acyclic noncontractible category of \claimref{B.7}, or the group
category of \claimref{B.10}, proves that the first inclusion is strict.
For the second inclusion use \claimref{F.8}. Since $F$ has no nontrivial
perfect subgroup, $BF$ is plus-local, and the two unit maps for $B\alpha$
are weak equivalences. Thus $P_+(B\alpha)$ is a weak equivalence if and
only if $B\alpha$ is, which it is not. But its homotopy cofiber is
integrally acyclic and connected and is therefore $\W_+$-aspherical.
The cofiber criterion of \claimref{B.5} proves the failure of Loc (6 b).
Finally $T(\W_+)=0$ follows from \claimref{F.6}.
\end{proof}

\begin{claim}{F.10}{Proved}{Nullifying an individual acyclic space does not supply reversible excision}
For every connected integrally acyclic small space $A$, let $P_A$ be
localization at $A\to *$, and let $\W_{P_A}$ be its class of equivalences
on categories. Then
\[
 \Winfty\subseteq\W_{P_A}\subseteq\W_+\subsetneq\WZ,
 \qquad T(\W_{P_A})=0.
\]
Every such $\W_{P_A}$ fails Loc (6 b). These assertions include
contractible $A$, for which $\W_{P_A}=\Winfty$.
\end{claim}
\begin{proof}
By \citet[Theorem 3.3]{FBerrickCasacuberta1999}, the unit
$X\to P_AX$ is a plus construction with respect to a particular perfect
normal subgroup of $\pi_1X$ and is an integral homology isomorphism.
The same theorem characterizes the local objects by the vanishing of the
subgroup swept out by maps from $A$. Any plus-local space has fundamental
group with no nontrivial perfect subgroup and is therefore $A$-null.
Equivalently, apply the construction to the universal perfect radical;
this gives the inclusion of equivalence classes
$\W_{P_A}\subseteq\W_+$. The other inclusions and $T=0$ follow from
\claimref{F.9} and \claimref{F.6}. If Loc (6 b) held, the minimality
theorem \claimref{B.6} would force $\WZ\subseteq\W_{P_A}$, contradicting
the strict upper inclusion. For contractible $A$, every space is $A$-null,
so the localization is the identity.
\end{proof}

\subsection{Retracts and the question on manuscript page 76}

\begin{claim}{F.11}{Proved}{The precise direct-factor condition}
The retract condition relevant to Loc (1 bis) is closure in the arrow
category, not merely closure under retracts of domains or codomains.
Explicitly, if $u:A\to B$ and $v:C\to D$ admit commuting squares
$(i_A,i_B):u\to v$ and $(r_A,r_B):v\to u$ with
$r_Ai_A=\id_A$ and $r_Bi_B=\id_B$, then $v\in\W$ must imply
$u\in\W$. Under the cofiber criterion of \claimref{B.5}, this condition
is equivalent to closure of the class of $\W$-aspherical categories under
retracts. For stable detection, closure of $\mathcal L$ under retracts
implies this condition for $R(\mathcal L)$. Conversely,
retract closure of $R(\mathcal L)$ alone does not recover all spectral
retract data of $\mathcal L$.
\end{claim}
\begin{proof}
Functoriality of the categorical mapping-cone model sends an arrow
retraction to a retraction of its cones. Thus closure of aspherical
objects under retracts, together with the cofiber criterion, gives
closure of arrows under retracts. Conversely, a retraction $A\to C\to A$
gives an arrow retraction of $A\to e$ from $C\to e$, so arrow retract
closure gives object retract closure without any cofiber hypothesis.
After applying $B$, $(-)_+$, and cofiber, an arrow retraction gives a
retraction $D(Bu)\to D(Bv)\to D(Bu)$ of spectra. This proves the
assertion for $R(\mathcal L)$. The last assertion follows already from
\claimref{F.7}: the entire stable kernel is not determined by the map
class. More precisely only its intersection with the suspension spectra
and the retractions induced there can be tested by this assignment.

For the modern fundamental localizers used in this paper, the arrow
retract condition is automatic:
\citet[Proposition 4.2.4, p.~183]{RootCisinski2006} proves strong saturation
from Definition 3.3.2, and a functor whose image is invertible is closed
under arrow retracts. This theorem uses the relative Theorem A; it is
not a license to omit Loc (1 bis) in a formal argument assuming only
the weaker absolute Loc (3) of the manuscript.
\end{proof}

\begin{claim}{F.12}{Proved}{Intersection and union aspherical imply both pieces aspherical}
Let $\W$ be a fundamental localizer in this paper's convention. Let
$U,V\subseteq X$ be full sieves with $U\cup V=X$, and put $K=U\cap V$.
If $K\to e$ and $X\to e$ belong to $\W$, then $U\to e$ and $V\to e$
belong to $\W$. Thus row XIIba-A-053 has an affirmative answer in the
specified convention. Full W(7 bis), and in particular its reverse half,
is unnecessary for this implication.
\end{claim}
\begin{proof}
If $K=\varnothing$, the hypothesis is $\varnothing\to e\in\W$.
Products preserve every fundamental localizer
\citep[Sorites 3.3.8(2)]{RootCisinski2006}. Taking products with an arbitrary
category $Y$ gives $\varnothing\to Y\in\W$; two out of three then
gives every functor in $\W$. The conclusion follows.

Suppose $K\ne\varnothing$, and choose an object $k\in K$. Since
$K\to e$ is in $\W$, its section $k:e\to K$ is in $\W$ by two out
of three. Let $I$ be the span category with three objects and arrows
$0\to1$, $0\to2$. Let
\[
 P=\int_I(U\leftarrow e\to V),\qquad
 Q=\int_I(U\leftarrow K\to V).
\]
The maps from $e$ select $k$ in each piece. The natural transformation
which is $k:e\to K$ at $0$ and the identity on the two arms induces a
functor $j:P\to Q$ over $I$. This functor is in $\W$. To verify the
last assertion directly from the available axioms, for a covariant
diagram $F:I\to\Cat$ the inclusion of $F(i)$ into
$(\int_I F)/i$ has the left adjoint which pushes
$(j,x,j\to i)$ forward to $(i,F(j\to i)x)$. It is therefore a weak
homotopy equivalence. Apply this observation to the two diagrams.
The induced functor on each comma category over $i$ is in $\W$ by two
out of three, because the map on $F(i)$ is. The relative Theorem A
gives $j\in\W$.

The functor $q:Q\to X$ is induced by the three inclusions into $X$.
It is a weak homotopy equivalence. One can check this without a
cofibration axiom: $Q/x$ is the Grothendieck construction of
$U/x\leftarrow K/x\to V/x$. If $x\in U\setminus V$, the sieve
condition on $U$ gives $V/x=K/x$, while $U/x$ has terminal object
$(x,\id_x)$; its homotopy pushout is therefore contractible. The case
$x\in V\setminus U$ is symmetric. If $x\in K$, all three comma
categories have terminal object $x$ and the span index is contractible.
The nerve comparison for Grothendieck constructions and Quillen's
Theorem A give the asserted weak equivalence. Thus $q\in\Winfty
\subseteq\W$, and $P\to e$ is in $\W$.

There are actual functors $s_U:U\to P$ and $r_U:P\to U$ with
$r_Us_U=\id_U$: $s_U$ is the arm inclusion; $r_U$ is the identity on
the $U$ arm, is constant at $k$ on the $V$ arm, and sends the central
object to $k$. On arrows from the central object to the $U$ arm it
uses the indicated arrow with source $k$; all arrows wholly on the
other arm, and its structural arrows, go to $\id_k$. These formulas
respect compositions. There are analogous $s_V,r_V$. Therefore
$U\to e$ and $V\to e$ are arrow retracts of $P\to e$, and
\claimref{F.11} proves the conclusion.
\end{proof}

\subsection{What the stable description does, and does not, classify}

\begin{claim}{F.13}{Proved}{The homology examples and the distinction from arbitrary localizing kernels}
For any spectrum $E$, define
$\mathcal A_E=\{V:E\wedge V=0\}$. Then
$R(\mathcal A_E)=\W_E$, where generalized homology of an unpointed
space $X$ means $\pi_*(E\wedge\Sigma^\infty_+X)$ in all integer
degrees. No ring structure on $E$ is required. Every full localizing
subcategory of $\Ho(\Sp)$ is a tensor ideal. Nevertheless the
definition of $R(\mathcal L)$ does not require a proof that
$\mathcal L=\mathcal A_E$ for some $E$.
\end{claim}
\begin{proof}
Smashing with $E$ preserves cofiber sequences. Thus
$E\wedge D(Bu)=0$ exactly when $u$ induces isomorphisms in all
$E$-homology groups. For the ideal assertion fix $V\in\mathcal L$.
The spectra $Z$ for which $V\wedge Z\in\mathcal L$ form a localizing
subcategory containing the sphere. The sphere generates spectra under
shifts, cofibers, and coproducts: attach sphere cells successively to
hit and then kill homotopy groups, and take the homotopy colimit of the
resulting cellular tower. Hence this subcategory is all of $\Ho(\Sp)$.
The construction and proof in \claimref{F.2} use only the localizing
closure properties, not representation as the acyclics of a smash
homology theory.
\end{proof}

\begin{claim}{F.14}{Open}{Two additional classification assertions not proved by the Galois connection}
We have not established either of the following assertions in this audit.
First, every accessible fundamental localizer satisfying Loc (6 b) is
fixed by $RT$. Equivalently, for such a $\W$ and every functor $g$,
\[
 D(Bg)\in\operatorname{loc}\{D(Bu):u\in\W\}
       \quad\Longrightarrow\quad g\in\W.
\]
Second, a specified arbitrary localizing subcategory of spectra need
not be declared a homological Bousfield class merely because it is
generated by a set. A concrete test not decided here is
\[
 \operatorname{loc}(H\mathbb F_p)
    \stackrel{?}{=}\{V:E\wedge V=0\}\quad\text{for some spectrum }E.
\]
We record the latter as an unresolved verification in this audit,
not as a new assertion about the present research status of that
particular example.
\end{claim}
The first missing implication is stronger than the algebraic description
of the image in \claimref{F.5}. The cofiber criterion proves invariance
under finite suspension and detects ordinary homotopy cofibers; it does
not, by itself, supply the comparison needed to transfer an arbitrary
stable construction, with desuspensions and infinite coproducts, back
to an unstable $\W$-equivalence. The unproved step is closure of the
unstable null class under all such stable generation relations. Even
for a set-generated reversible localizer, this implication must be
proved before replacing ``$RT$-fixed'' by ``reversible'' in the image
description. It holds for each of the concrete classes $\W_E$ already
in the image by \claimref{F.13} and for their intersections, by
\claimref{F.5}.

There is a useful exact comparison with the model-category literature.
\citet[Definition 2.1 and Proposition 2.3]{FBergerRatkovic2019}
call a left proper pointed model category \emph{excisive} when its
derived suspension is conservative, and identify this condition with
two-out-of-three for maps of cofibre sequences. Their Proposition 2.6
then detects homotopy pushouts by parallel homotopy cofibres.
However, the stable embedding in their Remark 2.9 assumes
\emph{fully faithful} derived suspension, called \emph{stably excisive},
as well as cofibrant generation. Conservativity alone does not meet
that hypothesis, so this theorem does not supply the missing
stable-detection implication above.

For the second question, \citet[Conjecture 9.1 and Proposition 9.2]{FHoveyPalmieri1999}
separate precisely the extra assertion that every localizing subcategory
of spectra is a homological Bousfield class. Their Proposition 9.2
reduces the general assertion to principal localizing subcategories;
it does not prove it. Set-generated localizing subcategories do admit
localization functors, as in
\citet[Theorem 3.7 and Proposition 3.8]{FCasacubertaGutierrezRosicky2014},
but that is a different statement. There are actual localizing
subcategories which are not homological Bousfield classes in the
$H\mathbb F_p$-\emph{local} stable category:
\citet[Proposition 6.4]{FWolcott2015} uses
\[
 \{V\text{ $H\mathbb F_p$-local}:
       [V,\Sigma^nM(p)]=0\text{ for all }n\in\mathbb Z\}.
\]
This is a proper nonzero localizing subcategory there, whereas that
category's Bousfield lattice has only two elements. Changing the
ambient category back to all spectra changes both coproducts and the
relevant tensor product; this example cannot be silently used as a
counterexample in $\Ho(\Sp)$. Our searches and their dates are listed
in the accompanying report. The definite failure of the bijection in
Lead 6 is \claimref{F.7}, independently of these two unresolved
classification assertions.

The second obligation can be made more precise without assuming the
general conjecture. Let $a(H\mathbb F_p)$ be a wedge of one representative
of each Bousfield class annihilating $H\mathbb F_p$; Ohkawa's set theorem
makes this a set-indexed wedge. Then
\[
 \operatorname{loc}(H\mathbb F_p)\subseteq
 \ker\bigl(a(H\mathbb F_p)\wedge-\bigr),
\]
and equality is necessary and sufficient for homological representability
of $\operatorname{loc}(H\mathbb F_p)$.
Indeed, the displayed kernel contains the generator and is contained
in every smash kernel containing that generator: any such kernel
is defined by a spectrum whose Bousfield class occurs in the wedge.
Thus this is the smallest possible homological kernel, rather than an
unspecified search for a representing spectrum. This is the specialized
argument of \citet[Proposition 9.2]{FHoveyPalmieri1999}.
Our additional check of \citet[Corollary 1.2]{FPuetzstueck2026}
concerns the distinct Hovey--Palmieri \emph{retract conjecture}
(their Conjecture 3.12), not their Conjecture 9.1 or the arrow-retract
axiom of \claimref{F.11}; it supplies no counterexample to the equality
just displayed.

